\documentclass{article}
\usepackage{amsmath}
\usepackage{algorithm}
\usepackage{algorithmic}
% A small synthetic paper used as IRP's offline golden example.
\newcommand{\R}{\mathbb{R}}

\title{Least Squares Recovers a Linear Trend in Ten Thousand Samples}
\author{IRP Examples}

\begin{document}
\maketitle

\begin{abstract}
We fit a line to noisy samples of $y = 2x + 1$ with ordinary least squares.
With $n = 10{,}000$ training samples the estimated slope is $2.00$ and the
test mean squared error is $0.0100$, matching the noise variance.
\end{abstract}

\section{Introduction}
\label{sec:intro}
Ordinary least squares (OLS) is the simplest estimator of a linear trend. We
report how close it gets to the true slope on synthetic data.

\section{Method}
\label{sec:method}
We model the data as
\begin{equation}
  y = w x + b + \varepsilon,
  \label{eq:model}
\end{equation}
and estimate $w$ and $b$ by minimizing the squared error
\begin{equation}
  \hat{w}, \hat{b} = \arg\min_{w, b} \sum_{i=1}^{n} (y_i - w x_i - b)^2,
  \label{eq:ols}
\end{equation}
whose closed form is
\begin{equation}
  \hat{w} = \frac{\sum_i (x_i - \bar{x})(y_i - \bar{y})}{\sum_i (x_i - \bar{x})^2},
  \qquad \hat{b} = \bar{y} - \hat{w}\bar{x}.
  \label{eq:closed}
\end{equation}

\begin{algorithm}[t]
\caption{Least squares fit}
\label{alg:fit}
\begin{algorithmic}
\STATE Compute the means $\bar{x}$ and $\bar{y}$
\STATE Compute $\hat{w}$ with Eq.~\ref{eq:closed}
\STATE Compute $\hat{b} = \bar{y} - \hat{w}\bar{x}$
\STATE Report the test mean squared error
\end{algorithmic}
\end{algorithm}


\section{Experiments}
\label{sec:experiments}
We draw $x \sim \mathcal{U}(0, 1)$ and $y = 2x + 1 + \varepsilon$ with
$\varepsilon \sim \mathcal{N}(0, 0.1^2)$, using 10{,}000 training and 10{,}000
test samples. Table~\ref{tab:main} reports the mean over three seeds.

\begin{table}[t]
\centering
\caption{Least squares on synthetic data, mean over three seeds.}
\label{tab:main}
\begin{tabular}{lcc}
\hline
Method & Slope & Test MSE \\
\hline
OLS (ours) & 2.00 & 0.0100 \\
\hline
\end{tabular}
\end{table}

The fit takes well under one second on a laptop CPU.

\appendix
\section{Implementation details}
\label{app:details}
Samples are drawn with a seeded pseudo-random generator.

\end{document}
